\section{Individuazione degli Stakeholder e dei Goal}

\subsection{Stakeholder e Goal}

Gli stakeholder principali del sistema sono l'\textbf{Operatore Sanitario} e il \textbf{Paziente Anziano}. I due attori usano la stessa piattaforma, ma con obiettivi e livelli di complessita' diversi. L'operatore deve gestire e interpretare i dati; il paziente deve svolgere azioni quotidiane semplici.

\subsubsection{Operatore Sanitario / Caregiver}

L'operatore sanitario e' il professionista che segue uno o piu' pazienti. Attraverso la dashboard puo' consultare l'elenco dei pazienti assegnati, entrare nel dettaglio di ciascun paziente, aggiornare alcune informazioni, configurare farmaci ed esercizi e monitorare i log giornalieri.

\begin{table}[H]
\centering
\footnotesize
\renewcommand{\arraystretch}{1.25}
\begin{tabular}{|p{3.2cm}|p{4.4cm}|p{4.7cm}|c|c|c|c|}
\hline
\textbf{Stakeholder} & \textbf{Goal} & \textbf{Subgoal} & \textbf{C} & \textbf{R} & \textbf{U} & \textbf{D} \\
\hline
Operatore Sanitario & Gestione pazienti & Visualizza lista pazienti & & X & & \\
\cline{3-7}
 & & Modifica dati del paziente & & & X & \\
\cline{3-7}
 & & Consulta dettaglio e metriche & & X & & \\
\hline
Operatore Sanitario & Gestione farmaci & Aggiunge farmaco & X & & & \\
\cline{3-7}
 & & Visualizza farmaci assegnati & & X & & \\
\cline{3-7}
 & & Elimina farmaco & & & & X \\
\hline
Operatore Sanitario & Gestione esercizi & Crea esercizio con step & X & & & \\
\cline{3-7}
 & & Visualizza esercizi assegnati & & X & & \\
\cline{3-7}
 & & Elimina esercizio & & & & X \\
\hline
Operatore Sanitario & Monitoraggio & Consulta umore, farmaci ed esercizi & & X & & \\
\cline{3-7}
 & & Richiede suggerimento AI & X & & & \\
\hline
\end{tabular}
\caption{Goal e operazioni CRUD dell'Operatore Sanitario}
\label{tab:goal-operatore}
\end{table}

\subsubsection{Paziente Anziano}

Il paziente anziano usa l'applicazione per comunicare poche informazioni essenziali. Non gestisce il proprio piano terapeutico, ma interagisce con quanto configurato dall'operatore: vede i farmaci, conferma l'assunzione, vede gli esercizi, li completa e registra l'umore del giorno.

\begin{table}[H]
\centering
\footnotesize
\renewcommand{\arraystretch}{1.25}
\begin{tabular}{|p{3.2cm}|p{4.2cm}|p{4.9cm}|c|c|c|c|}
\hline
\textbf{Stakeholder} & \textbf{Goal} & \textbf{Subgoal} & \textbf{C} & \textbf{R} & \textbf{U} & \textbf{D} \\
\hline
Paziente Anziano & Accesso & Login e registrazione associata a un operatore & X & X & & \\
\hline
Paziente Anziano & Umore & Registra o aggiorna umore quotidiano & X & X & X & \\
\hline
Paziente Anziano & Esercizi fisici & Visualizza lista esercizi & & X & & \\
\cline{3-7}
 & & Completa esercizio guidato & X & X & & \\
\hline
Paziente Anziano & Farmaci & Visualizza farmaci per fascia oraria & & X & & \\
\cline{3-7}
 & & Conferma o annulla assunzione & X & X & X & \\
\hline
\end{tabular}
\caption{Goal e operazioni CRUD del Paziente Anziano}
\label{tab:goal-paziente}
\end{table}

\newpage
\subsection{Diagramma di Contesto}

Il diagramma di contesto rappresenta A.L.I.C.E. come una scatola nera posta al centro dell'interazione tra utenti e servizi esterni. L'operatore sanitario invia al sistema configurazioni e aggiornamenti relativi a pazienti, farmaci ed esercizi; in cambio riceve dashboard, metriche, storico dei log e suggerimenti sintetici. Il paziente invia conferme giornaliere, completamenti degli esercizi e valore dell'umore; in cambio riceve una home semplificata, l'elenco delle attivita' e le istruzioni operative.

Sul lato dei servizi, Supabase fornisce autenticazione, persistenza dei dati e controllo di visibilita' tramite Row Level Security. La Groq API viene usata solo dalla route server-side dedicata al suggerimento AI: il browser non riceve la chiave del servizio esterno e il testo generato e' destinato all'operatore.

\begin{figure}[H]
\centering
\begin{tikzpicture}[
    entity/.style={rectangle, draw=aliceblue, fill=alicelight, thick, minimum width=3.5cm, minimum height=1cm, align=center, font=\small\bfseries},
    system/.style={rectangle, draw=aliceblue, fill=aliceblue!15, thick, rounded corners=6pt, minimum width=4.5cm, minimum height=2cm, align=center, font=\large\bfseries},
    service/.style={rectangle, draw=alicegreen, fill=green!8, thick, rounded corners=4pt, minimum width=3.6cm, minimum height=1cm, align=center, font=\small\bfseries},
    flow/.style={-{Stealth[length=2.5mm]}, thick, color=aliceblue},
    flowlabel/.style={font=\scriptsize, fill=white, inner sep=2pt, align=center},
]
\node[system] (alice) {A.L.I.C.E.\\Applicazione Web};
\node[entity, above left=2.4cm and 2.8cm of alice] (operatore) {Operatore\\Sanitario};
\node[entity, above right=2.4cm and 2.8cm of alice] (paziente) {Paziente\\Anziano};
\node[service, below left=2.2cm and 1.8cm of alice] (supabase) {Supabase\\Auth + PostgreSQL};
\node[service, below right=2.2cm and 1.8cm of alice] (groq) {Groq API\\AI};

\draw[flow] (operatore) -- node[flowlabel] {configura pazienti,\\farmaci, esercizi} (alice);
\draw[flow] (alice) -- node[flowlabel] {dashboard, metriche,\\storici, suggerimento} (operatore);
\draw[flow] (paziente) -- node[flowlabel] {umore, conferme,\\completamenti} (alice);
\draw[flow] (alice) -- node[flowlabel] {home, farmaci,\\esercizi guidati} (paziente);
\draw[flow] (alice) -- node[flowlabel] {sessioni e query} (supabase);
\draw[flow] (supabase) -- node[flowlabel] {dati e RLS} (alice);
\draw[flow] (alice) -- node[flowlabel] {dati sintetici} (groq);
\draw[flow] (groq) -- node[flowlabel] {frase operativa} (alice);
\end{tikzpicture}
\caption{Diagramma di contesto del sistema A.L.I.C.E.}
\label{fig:contesto}
\end{figure}

\subsection{Scambio di Valore}

\begin{figure}[H]
\centering
\begin{tikzpicture}[
    system/.style={circle, draw=aliceblue, fill=aliceblue!12, thick, minimum size=2.7cm, align=center, font=\bfseries},
    ai/.style={rectangle, draw=aliceorange, fill=aliceorange!10, thick, rounded corners=4pt, minimum width=2.8cm, minimum height=0.8cm, align=center, font=\scriptsize\bfseries},
    flow/.style={-{Stealth[length=2.5mm]}, thick},
    label/.style={font=\scriptsize, fill=white, inner sep=2pt, align=center},
]

\node[system] (alice) at (0, 0) {A.L.I.C.E.};
\node[ai] (ai) at (0, -2.3) {Feature AI\\Suggerimento};

% omino operatore
\coordinate (op) at (-5.5, 0);
\draw[aliceblue, thick] ($(op)+(0,0.8)$) circle (0.35);
\draw[aliceblue, thick] ($(op)+(0,0.45)$) -- ($(op)+(0,-0.55)$);
\draw[aliceblue, thick] ($(op)+(-0.65,0.1)$) -- ($(op)+(0.65,0.1)$);
\draw[aliceblue, thick] ($(op)+(0,-0.55)$) -- ($(op)+(-0.55,-1.25)$);
\draw[aliceblue, thick] ($(op)+(0,-0.55)$) -- ($(op)+(0.55,-1.25)$);
\node[align=center, font=\small\bfseries] at ($(op)+(0,-1.75)$) {Operatore\\Sanitario};

% omino paziente
\coordinate (pz) at (5.5, 0);
\draw[alicegreen, thick] ($(pz)+(0,0.8)$) circle (0.35);
\draw[alicegreen, thick] ($(pz)+(0,0.45)$) -- ($(pz)+(0,-0.55)$);
\draw[alicegreen, thick] ($(pz)+(-0.65,0.1)$) -- ($(pz)+(0.65,0.1)$);
\draw[alicegreen, thick] ($(pz)+(0,-0.55)$) -- ($(pz)+(-0.55,-1.25)$);
\draw[alicegreen, thick] ($(pz)+(0,-0.55)$) -- ($(pz)+(0.55,-1.25)$);
\node[align=center, font=\small\bfseries] at ($(pz)+(0,-1.75)$) {Paziente\\Anziano};

\draw[flow, color=aliceblue] (-4.5, 0.7) -- node[label] {pazienti, farmaci,\\esercizi} (-1.45, 0.75);
\draw[flow, color=aliceblue] (-1.45, -0.7) -- node[label] {metriche, storico,\\supporto decisionale} (-4.5, -0.7);
\draw[flow, color=alicegreen] (4.5, 0.7) -- node[label] {umore, conferme,\\completamenti} (1.45, 0.75);
\draw[flow, color=alicegreen] (1.45, -0.7) -- node[label] {interfaccia semplice,\\routine giornaliera} (4.5, -0.7);
\draw[flow, color=aliceorange] (alice) -- node[label, right] {genera frase\\per operatore} (ai);
\end{tikzpicture}
\caption{Scambio di valore con Operatore, Paziente e feature AI}
\label{fig:scambio-valore}
\end{figure}

\newpage
